\documentclass[11pt]{article}

\usepackage[margin=1in]{geometry}
\usepackage{amsmath, amssymb, amsthm,cancel}
\usepackage{mathtools}
\usepackage{enumitem}
\usepackage{graphicx}
\usepackage{booktabs}
\usepackage{float}
\usepackage{xcolor}
\usepackage{listings}

\newcommand{\R}{\mathbb{R}}
\newcommand{\norm}[1]{\left\lVert #1\right\rVert}
\newcommand{\answerline}{\noindent\rule{\textwidth}{0.4pt}}

\lstset{
    basicstyle=\ttfamily\small,
    breaklines=true,
    frame=single,
    captionpos=b
}

\title{Math 609 -- Homework 5}
\author{Darin Peries}
\date{\today}

\begin{document}

\maketitle

\section*{Problem 1}

\subsection*{(a) Newton's method for roots}

Applying Newton's method to solve
\[
x^2=a
\]
gives the square-root iteration
\[
t_{n+1}=\frac{1}{2}\left(t_n+\frac{a}{t_n}\right).
\]

\subsubsection*{Answer:}

\subsubsection*{Derivation of the square root formula}

\noindent In order to get Newton's method to compute $\sqrt{a}$, we need to transform the function $x^2=a$ so that $x=\sqrt{a}$.
when $f(x)=0$. The function with the necessary transformation is

$$f(x) = x^2 - a$$

\noindent The derivative of this funtion is 

$$f'(x) = 2x$$

\noindent and Newton's method gives

\[t_{n+1} = t_n - \frac{f(t_n)}{f'(t_n)} = t_n - \frac{t_n^2-a}{2t_n} =
 t_n - \frac{t_n^{\cancel{2}}}{2\cancel{t_n}} + \frac{a}{2t_n} = \frac{t_n}{2} + \frac{a}{2t_n}=
 \frac{1}{2}(t_n + \frac{a}{t_n})\]


\subsubsection*{Cubed quartic root derivation}
\noindent Following the derivation of Newton's method for the square root, the same can be done to calculate the cubed root.
First, transform $x^3$ to a function, where $x = \sqrt[3]{a}$ when $f(x) = 0$.
The transformed function and following derivation is

$$f(x) = x^3 - a$$

$$f'(x) = 3x^2 $$

\[ t_{n+1} = t_n - \frac{t_n^3-a}{3t_n^2} = t_n - \frac{t_n^{\cancel{3}}}{3\cancel{t_n^2}} + \frac{a}{3t_n^2}
 = \frac{2t_n}{3} + \frac{a}{3t_n^2} = \frac{1}{3}(2t_n + \frac{a}{t_n^2})\]

\noindent Using the prior derivations, the algorithm/formula to get the subsequent roots is

$$\frac{1}{root}((root-1)t_n + \frac{a}{t_n^{root-1}})$$

\noindent For instance, the quartic root is

\[ t_{n+1} = t_n - \frac{t_n^4-a}{4t_n^3} = \frac{3t_n}{4} + \frac{a}{4t_n^3} = \frac{1}{4}(3t_n + \frac{a}{t_n^3}) \]

\subsubsection*{a=5}

\noindent To get the squared root of $5$, the initial guess, $t_0$ needs to be sufficietly close
to the root of $f(x)$. For this case, $t_0=2$ would be the best choice, as it is the closest
square root to $\sqrt{5}$.\\

\noindent After 4 iterations, the approximate squared root of 5 is $2.236$.

\subsection*{(b) Fixed-point iteration}

Consider
\[
f(x)=x^3-x+1.
\]

\noindent Find an appropriate interval and a function \(g\) such that the fixed-point iteration
\[
x_{n+1}=g(x_n)
\]
generates a sequence converging to a root of \(f\).

\subsubsection*{Answer:}

\subsubsection*{Finding the interval}
To get the proper interval, a chosen function $g(x)$ needs to map back to the interval, $g([a,b])\subseteq [a,b]$.
If we choose a function $g(x_n)=\sqrt[3]{x-1}$, a good interval to start is $[-2, 0]$, as $g(-2)=\sqrt[3]{-3}$ and $g(0)=-1$, 
thus the function maps into itself.

\subsubsection*{Show $g(x_n)$ is a contraction}
In order to show that $g(x)$ converges on the interval $[-2, 0]$ the function needs to be a contraction map.
This is shown by

$$|g(x) - g(y)| \leq \lambda |x - y|$$

\noindent where $\lambda < 1$ and positive. Using the Mean Value Theorem, $\lambda$ is the 
absolute value of the derivative of a point, inside the interval, if the absolute value of the derivative is
less than one.

$$g'(x) = \frac{1}{3}(x-1)^{-2/3}$$

\noindent and

$$\lambda = |g'(0)| = |\frac{1}{3\sqrt[3]{4}}| < 1$$

\noindent Because $g(x)$ maps to itself and is a contraction on the interval, the iteration method will
produce a sequence of solutions that converge. After $13$ iterations, with an initial guess of $0$ and stopping criteria of $1\times 10^{-9}$, the approximate solution is $-1.325$.


\section*{Problem 2}

Let \(\alpha\) be a solution of
\[
f(x)=0,
\]
and let \(\{t_n\}\) be the sequence generated by Newton's method:
\[
t_{n+1}=t_n-\frac{f(t_n)}{f'(t_n)}.
\]

\noindent Show that
\[
\alpha-t_{n+1}
=
-\frac{1}{2}\frac{f''(\xi_n)}{f'(t_n)}
(\alpha-t_n)^2,
\]
where \(\xi_n\) lies between \(t_n\) and \(\alpha\).

\bigskip

\noindent\textbf{Answer:}\\

\noindent By using the taylor series expansion, the function 
$f(t_n)$ with $\alpha$ as the center, which gives

\[ f(\alpha) = f(t_n) + f'(t_n)(\alpha-t_n) + \frac{f''(\xi_n)}{2}(\alpha-t_n)^2 \]
\[ 0 = f(t_n) + f'(t_n)(\alpha-t_n) + \frac{f''(\xi_n)}{2}(\alpha-t_n)^2 \]
\[ f(t_n) = -f'(t_n)(\alpha-t_n) - \frac{f''(\xi_n)}{2}(\alpha-t_n)^2 \]
\noindent Newton's iteration can then be expanded

\[ 
  t_{n+1} = t_n - \frac{1}{f'(t_n)}(-f'(t_n)(\alpha-t_n) - \frac{f''(\xi_n)}{2}(\alpha-t_n)^2)
  =t_n + \frac{f'(t_n)}{f'(t_n)}(\alpha - t_n) + \frac{f''(\xi)}{2f'(t_n)}(\alpha - t_n)^2
\]

\[
  t_{n+1}=t_n + (\alpha - t_n) + \frac{f''(\xi)}{f'(2t_n)}(\alpha - t_n)^2
\]

\[
  t_{n+1} - \alpha = \frac{1}{2}\frac{f''(\xi)}{f'(t_n)}(\alpha - t_n)^2  
\]

\[
  \alpha - t_{n+1} = -\frac{1}{2}\frac{f''(\xi)}{f'(t_n)}(\alpha-t_n)^2
\]

\vspace{6in}

\newpage

\section*{Problem 3}

The nonlinear system
\[
x_1^2-10x_1+x_2^2+8=0,
\]
\[
x_1x_2^2+x_1-10x_2+8=0
\]
can be transformed into the fixed-point system
\[
x_1=g_1(x_1,x_2)=\frac{x_1^2+x_2^2+8}{10},
\]
\[
x_2=g_2(x_1,x_2)=\frac{x_1x_2^2+x_1+8}{10}.
\]

\subsection*{(a) show that $G=(g_1,g_2)^T$ has a unique fixed point}

The interval for $x_1,x_2$ is $[0,1.5]$. First, $G$ needs to map to itself, 
$G:[0,1.5] \times [0,1.5] \subseteq [0,1.5] \times [0,1.5]$.

\[
\begin{aligned}
  &g_1(0, 0) =  \frac{4}{5}\\
  &g_1(1.5, 1.5) = \frac{2.25 + 2.25 + 8}{10} = \frac{12.5}{10} = 1.25 \\
  &g_2(0, 0) = \frac{4}{5}\\
  &g_2(1.5, 1.5) = \frac{3.375 + 1.5 + 8}{10} = \frac{12.875}{10} = 1.2875
\end{aligned}
\]

\noindent From evaluating the functions at the ends of the interval, it is evident that
the function maps the interval back into itself. Next, the function $G$ needs to also be a 
contraction. The jacobian of $G$ is 

\[J_G = 
\begin{pmatrix}
    \frac{x_1}{5} & \frac{x_2}{5} \\
    \frac{x_2^2+1}{10} & \frac{x_1x_2}{5} \\
\end{pmatrix}
\]

\noindent In order for $G$ to be a contraction, the maximum of $\|J_G\|_{\infty}$ needs to be
less than one. Because of the increasing monoticity of the function, the maximum is at $x_1=1.5$ and $x_2=1.5$.

\[
\begin{aligned}
  &J_G(1.5, 1.5) = 
  \begin{pmatrix}
    0.3 & 0.3 \\
    0.325 & 0.45
  \end{pmatrix} \\
  &\|J_G\|_{\infty} = \max (|0.3 + 0.3|, |0.325 + 0.45|) = 0.775 \\
  & 0.775 < 1
\end{aligned}
\]

\noindent Because $\|J_G\|_{\infty} < 1$, $G$ is a contraction. Due to the 
Banach fixed point theorem, $G$ has a single fixed point.

\subsection*{(b) Result from code.}
After running for $22$ iterations with a starting point of $(1.5, 1.5)^T$ and a stopping tolerance for $1 \times 10^{-9}$, the result of the Fixed-Point iteration code is $(1.0, 1.0)^T$.

\newpage

\section*{Problem 4}

\subsection*{(a) First derivative formula}

Show that
\[
f'(t)\approx
\frac{1}{2h}
\left[
-3f(t)+4f(t+h)-f(t+2h)
\right].
\]

\noindent Determine the order of the truncation error.

\bigskip

\noindent\textbf{Derivation:}

\noindent By using the Taylor expansion for each point in the approximation,

\[
  \begin{aligned}
    &f(t+h) = f(t) + f'(t)h + \frac{f''(t)h^2}{2} + \frac{f'''(t)h^3}{6} + \frac{f''''(t)h^4}{24}\\
    &f(t) = f(t)\\
    &f(t+2h) = f(t) + f'(t)2h + 2f''(t)h^2 + \frac{4f'''(t)h^3}{3} + \frac{2f''''(t)h^4}{3}
  \end{aligned}
\]

\noindent By substituting the expansions in the approximation,

\[
\begin{aligned}
  &-3f(t) + 4f(t+h) = f(t) + 4f'(t)h + 2f''(t)h^2 + \frac{2f'''(t)h^3}{3} + \frac{f''''(t)h^4}{6}\\
  &\frac{1}{2h}(-3f(t) + 4f(t+h)) - f(t+2h) = \frac{1}{2h}(2f'(t)h - \frac{2h^3f'''(\xi)}{3})=f'(t) - \frac{h^2f'''(\xi)}{3}\\
  &\frac{1}{2h}[-3f(t) + 4f(t+h) - f(t+2h)] -f'(t) = -\frac{h^2f'''(\xi)}{3} 
\end{aligned}
\]

\noindent The error term is of order $O(h^2)$.

\subsection*{(b) Second derivative formula}

Show that
\[
f''(t)\approx
\frac{1}{12h^2}
\left[
-f(t+2h)+16f(t+h)-30f(t)+16f(t-h)-f(t-2h)
\right].
\]

Determine the order of the truncation error.

\bigskip

\noindent\textbf{Derivation:}

\noindent Using the logic from the prior part, expanding all the points gives,

\[
\begin{aligned}
  &f(t+h) = f(t) + f'(t)h + \frac{f''(t)h^2}{2} + \frac{f'''(t)h^3}{6} + \frac{f''''(t)h^4}{24} \\
  &+ \frac{f^{(5)}(t)h^5}{120} + \frac{f^{(6)}(t)h^6}{720}\\
  &f(t) = f(t)\\
  &f(t+2h) = f(t) + f'(t)2h + 2f''(t)h^2 + \frac{4f'''(t)h^3}{3} + \frac{2f''''(t)h^4}{3} \\
  &+ \frac{4f^{(5)}(t)h^5}{15} + \frac{4f^{(6)}(t)h^6}{45}\\
  &f(t-h) = f(t) - f'(t)h + \frac{f''(t)h^2}{2} - \frac{f'''(t)h^3}{6} + \frac{f''''(t)h^4}{24}\\
  &- \frac{f^{(5)}(t)h^5}{120} + \frac{f^{(6)}(t)h^6}{720}\\
  &f(t-2h) = f(t) - f'(t)2h + 2f''(t)h^2 - \frac{4f'''(t)h^3}{3} + \frac{2f''''(t)h^4}{3}\\
  &- \frac{4f^{(5)}(t)h^5}{15} + \frac{4f^{(6)}(t)h^6}{45}\\
\end{aligned}
\]

\noindent substituting into the approximations

\[
\begin{aligned}
  &-f(t+2h) + 16f(t+h) = 15f(t) + 14f'(t)h + 6f''(t)h^2 + \frac{4 f'''(t)h^3}{3}\\
  & -\frac{2h^5f^{(5)}(t)}{15} - \frac{h^6f^{(6)}(t)}{15}\\
  &-30f(t) + 16f(t-h) = -14f(t) - 16f'(t)h + 8f''(t)h^2 - \frac{8f'''(t)h^3}{3} + \frac{2f''''(t)h^4}{3}\\
  &- \frac{2h^5f^{(5)}(t)}{15} + \frac{h^6f^{(6)}(t)}{45}\\
  &(-f(t+2h) + 16f(t+h)) + (-30f(t) + 16f(t-h)) = f(t) -2f'(t)h + 14f''(t)h^2  - \frac{4f'''(t)h^3}{3} + \frac{2f''''(t)h^4}{3}\\
  & - \frac{4h^5f^{(5)}(t)}{15}\\
  &((-f(t+2h) + 16f(t+h)) + (-30f(t) + 16f(t-h))) -f(t-2h) = 12f''(t)h^2 - \frac{6h^6f^{(6)}(t)}{45} \\
  &\frac{1}{12h^2}(-f(t+2h) + 16f(t+h) - 30f(t) + 16f(t-h) -f(t-2h)) = f''(t) - \frac{h^4f^{(6)}(t)}{90}\\
  &\frac{1}{12h^2}(-f(t+2h) + 16f(t+h) - 30f(t) + 16f(t-h) -f(t-2h)) - f''(t) = - \frac{h^4f^{(6)}(t)}{90}
\end{aligned}
\]

\noindent The error is of order $O(h^4)$.

\section*{Problem 5}

Consider the two-point boundary-value problem
\[
-\phi''+(4x^2+2)\phi=0,
\qquad
\phi(0)=1,
\qquad
\phi(1)=e,
\]
whose exact solution is
\[
\phi(x)=e^{x^2}.
\]

Write a program to solve the problem using the second-order finite-difference scheme. Use step sizes
\[
h=2^{-p},
\qquad
p=1,2,\ldots,14.
\]

Solve the resulting linear system using:
\begin{enumerate}
    \item LU factorization / Thomas algorithm;
    \item Gradient descent;
    \item Conjugate gradient.
\end{enumerate}

For the iterative methods, state the initial guess and stopping criterion. Include your code, plots, and the requested error tables.\\\\

\noindent All iterative methods were ran with $x_0=0$ and stopping tolerance of $1\times10^{-9}$

\subsection*{Method 1: Thomas algorithm}


\section*{Error Tables}

% Thomas algorithm results
\begin{table}[htbp]
\centering
\caption{Errors using the Thomas algorithm.}
\resizebox{\textwidth}{!}{%
\begin{tabular}{ccccc}
\toprule
$h$ &
$\lVert u_h-\phi_h\rVert_\infty$ &
$\lVert u_h-\phi_h\rVert_\infty/h$ &
$\lVert u_h-\phi_h\rVert_\infty/h^2$ &
$\lVert u_h-\phi_h\rVert_\infty/h^3$ \\
\midrule
$5.0000\times10^{-1}$ & $6.8077\times10^{-2}$ & $1.3615\times10^{-1}$ & $2.7231\times10^{-1}$ & $5.4462\times10^{-1}$ \\
$2.5000\times10^{-1}$ & $2.0012\times10^{-2}$ & $8.0048\times10^{-2}$ & $3.2019\times10^{-1}$ & $1.2808$ \\
$1.2500\times10^{-1}$ & $5.4792\times10^{-3}$ & $4.3834\times10^{-2}$ & $3.5067\times10^{-1}$ & $2.8054$ \\
$6.2500\times10^{-2}$ & $1.3851\times10^{-3}$ & $2.2161\times10^{-2}$ & $3.5458\times10^{-1}$ & $5.6732$ \\
$3.1250\times10^{-2}$ & $3.4860\times10^{-4}$ & $1.1155\times10^{-2}$ & $3.5697\times10^{-1}$ & $11.4230$ \\
$1.5625\times10^{-2}$ & $8.7213\times10^{-5}$ & $5.5816\times10^{-3}$ & $3.5722\times10^{-1}$ & $22.8624$ \\
$7.8125\times10^{-3}$ & $2.1807\times10^{-5}$ & $2.7913\times10^{-3}$ & $3.5729\times10^{-1}$ & $45.7330$ \\
$3.9063\times10^{-3}$ & $5.4521\times10^{-6}$ & $1.3957\times10^{-3}$ & $3.5731\times10^{-1}$ & $91.4702$ \\
$1.9531\times10^{-3}$ & $1.3630\times10^{-6}$ & $6.9787\times10^{-4}$ & $3.5731\times10^{-1}$ & $182.9429$ \\
$9.7656\times10^{-4}$ & $3.4076\times10^{-7}$ & $3.4894\times10^{-4}$ & $3.5731\times10^{-1}$ & $365.8881$ \\
$4.8828\times10^{-4}$ & $8.5188\times10^{-8}$ & $1.7447\times10^{-4}$ & $3.5731\times10^{-1}$ & $731.7615$ \\
$2.4414\times10^{-4}$ & $2.1301\times10^{-8}$ & $8.7249\times10^{-5}$ & $3.5737\times10^{-1}$ & $1463.7983$ \\
$1.2207\times10^{-4}$ & $5.3250\times10^{-9}$ & $4.3622\times10^{-5}$ & $3.5735\times10^{-1}$ & $2927.4504$ \\
$6.1035\times10^{-5}$ & $4.1937\times10^{-9}$ & $6.8710\times10^{-5}$ & $1.1257$ & $18444.1299$ \\
\bottomrule
\end{tabular}%
}
\end{table}

% Conjugate Gradient results
\begin{table}[htbp]
\centering
\caption{Errors using the Conjugate Gradient method.}
\resizebox{\textwidth}{!}{%
\begin{tabular}{ccccc}
\toprule
$h$ &
$\lVert u_h-\phi_h\rVert_\infty$ &
$\lVert u_h-\phi_h\rVert_\infty/h$ &
$\lVert u_h-\phi_h\rVert_\infty/h^2$ &
$\lVert u_h-\phi_h\rVert_\infty/h^3$ \\
\midrule
$5.0000\times10^{-1}$ & $6.8077\times10^{-2}$ & $1.3615\times10^{-1}$ & $2.7231\times10^{-1}$ & $5.4462\times10^{-1}$ \\
$2.5000\times10^{-1}$ & $2.0012\times10^{-2}$ & $8.0048\times10^{-2}$ & $3.2019\times10^{-1}$ & $1.2808$ \\
$1.2500\times10^{-1}$ & $5.4792\times10^{-3}$ & $4.3834\times10^{-2}$ & $3.5067\times10^{-1}$ & $2.8054$ \\
$6.2500\times10^{-2}$ & $1.3851\times10^{-3}$ & $2.2161\times10^{-2}$ & $3.5458\times10^{-1}$ & $5.6732$ \\
$3.1250\times10^{-2}$ & $3.4860\times10^{-4}$ & $1.1155\times10^{-2}$ & $3.5697\times10^{-1}$ & $11.4230$ \\
$1.5625\times10^{-2}$ & $8.7213\times10^{-5}$ & $5.5816\times10^{-3}$ & $3.5722\times10^{-1}$ & $22.8624$ \\
$7.8125\times10^{-3}$ & $2.1807\times10^{-5}$ & $2.7913\times10^{-3}$ & $3.5729\times10^{-1}$ & $45.7330$ \\
$3.9063\times10^{-3}$ & $5.4521\times10^{-6}$ & $1.3957\times10^{-3}$ & $3.5731\times10^{-1}$ & $91.4702$ \\
$1.9531\times10^{-3}$ & $1.3630\times10^{-6}$ & $6.9787\times10^{-4}$ & $3.5731\times10^{-1}$ & $182.9429$ \\
$9.7656\times10^{-4}$ & $3.4076\times10^{-7}$ & $3.4894\times10^{-4}$ & $3.5731\times10^{-1}$ & $365.8881$ \\
$4.8828\times10^{-4}$ & $8.5190\times10^{-8}$ & $1.7447\times10^{-4}$ & $3.5731\times10^{-1}$ & $731.7777$ \\
$2.4414\times10^{-4}$ & $2.1297\times10^{-8}$ & $8.7233\times10^{-5}$ & $3.5731\times10^{-1}$ & $1463.5295$ \\
$1.2207\times10^{-4}$ & $5.3239\times10^{-9}$ & $4.3613\times10^{-5}$ & $3.5728\times10^{-1}$ & $2926.8251$ \\
$6.1035\times10^{-5}$ & $4.1742\times10^{-9}$ & $6.8390\times10^{-5}$ & $1.1205$ & $18358.4111$ \\
\bottomrule
\end{tabular}%
}
\end{table}

% Gradient Descent results
\begin{table}[htbp]
\centering
\caption{Errors using the Gradient Descent method.}
\resizebox{\textwidth}{!}{%
\begin{tabular}{ccccc}
\toprule
$h$ &
$\lVert u_h-\phi_h\rVert_\infty$ &
$\lVert u_h-\phi_h\rVert_\infty/h$ &
$\lVert u_h-\phi_h\rVert_\infty/h^2$ &
$\lVert u_h-\phi_h\rVert_\infty/h^3$ \\
\midrule
$5.0000\times10^{-1}$ & $6.8077\times10^{-2}$ & $1.3615\times10^{-1}$ & $2.7231\times10^{-1}$ & $5.4462\times10^{-1}$ \\
$2.5000\times10^{-1}$ & $2.0012\times10^{-2}$ & $8.0048\times10^{-2}$ & $3.2019\times10^{-1}$ & $1.2808$ \\
$1.2500\times10^{-1}$ & $5.4247\times10^{-3}$ & $4.3398\times10^{-2}$ & $3.4718\times10^{-1}$ & $2.7775$ \\
$6.2500\times10^{-2}$ & $1.3579\times10^{-1}$ & $2.1726$ & $34.7624$ & $556.1979$ \\
$3.1250\times10^{-2}$ & $9.2697\times10^{-1}$ & $29.6631$ & $949.2195$ & $30375.0233$ \\
$1.5625\times10^{-2}$ & $1.4913$ & $95.4429$ & $6108.3461$ & $390934.1507$ \\
$7.8125\times10^{-3}$ & $1.9168$ & $245.3517$ & $31405.0154$ & $4019841.9717$ \\
$3.9063\times10^{-3}$ & $2.2328$ & $571.5924$ & $146327.6659$ & $37459882.4808$ \\
$1.9531\times10^{-3}$ & $2.4386$ & $1248.5529$ & $639259.0594$ & $327300638.4105$ \\
$9.7656\times10^{-4}$ & $2.5624$ & $2623.9028$ & $2686876.4235$ & $2751361457.6764$ \\
$4.8828\times10^{-4}$ & $2.6333$ & $5393.0959$ & $11045060.4320$ & $22620283764.6882$ \\
$2.4414\times10^{-4}$ & $2.6727$ & $10947.5173$ & $44841030.7701$ & $183668862034.2852$ \\
$1.2207\times10^{-4}$ & $2.6941$ & $22070.3813$ & $180800563.6421$ & $1481118217355.7407$ \\
$6.1035\times10^{-5}$ & $2.7056$ & $44328.4526$ & $726277367.5724$ & $11899328390305.924$ \\
\bottomrule
\end{tabular}%
}
\end{table}

\newpage

\section*{Plots}

\begin{figure}[htbp]
  \centering
  \includegraphics[width=0.6\textwidth]{thomas_errors.png}
  \caption{Error quantities for the Thomas algorithm.}
  \label{fig:thomas-errors}
\end{figure}

\begin{figure}[htbp]
  \centering
    \includegraphics[width=0.6\textwidth]{gd_errors.png}
    \caption{Error quantities for the Gradient Descent algorithm.}
    \label{fig:Descent-errors}
\end{figure}

\begin{figure}[htbp]
  \centering
    \includegraphics[width=0.6\textwidth]{cg_errors.png}
    \caption{Error quantities for the Conjugate Gradient algorithm.}
    \label{fig:Conjugate-errors}
\end{figure}

\end{document}